\section{An exact algorithm for the observable chart}
\label{sec:v18-chart}
The chart in Lemma~\ref{lem:v17-chart} can be found without taking its
dimension as an oracle. We state the input and output explicitly. Rational
polynomials are encoded by finite lists of monomial exponent vectors and
rational coefficients in reduced integer form. A product of simplexes is
encoded by the number of entries in each row. Eliminating the last entry
of every nontrivial row leaves a full-dimensional rational polytope
$\Omega\subset\R^r$. Zero-dimensional factors are omitted.

\begin{proposition}[Effective observable quotient and pullback]
\label{prop:v18-chart}
The input consists of the preceding simplex description, a rational
polynomial map $b:\Omega\to\R^N$, and finitely many rational affine
functionals $f_i$ on $\R^N$. There is a terminating algorithm which returns
its observable affine dimension $a$, a rational basis and coordinate map
for the observable quotient, and, when $a>0$, rational parameter points
$q_0,q_1,\ldots,q_a$ with the following properties. With
\[
 z(q)=(f_i(b(q))-f_i(b(q_0)))_i,\qquad v_j=z(q_j),
\]
the vectors $v_j$ are a basis of the observable difference space, and
all coordinates $t(q)$ in that basis satisfy $|t_j(q)|<2$. The output also
contains the rational polynomial cube map $x(q)=(t(q)+2)/4$, a rational
upper bound for the observable diameter, and the exact pullback of any
subsequently supplied rational polynomial in $x$.
The algorithm outputs the universal determinant formula it has decided.
It makes no polynomial-time or coefficient-height claim.
\end{proposition}
\begin{proof}
Choose a rational relative interior point $q_0$ of the input simplexes.
Expand the vector polynomial $z(q)$ in the free row coordinates. Let $V$
be the rational span of all its coefficient vectors, including its constant
coefficient. This equals $\Span z(\Omega)$: a linear functional vanishes
on $z(\Omega)$ exactly when its polynomial vanishes on the nonempty
interior of $\Omega$, hence identically, hence on every coefficient.
Gaussian elimination therefore gives $a=\dim V$ and a rational basis
matrix $U$. Choosing independent rows gives a rational left inverse $L$
with $LU=I_a$. Thus $y(q)=Lz(q)$ provides exact quotient coordinates.
If $a=0$, all discrepancies are constant and no chart is required.

For $a>0$ let
\[
 D_0=\max_{u_1,\ldots,u_a\in\Omega}
             |\det(y(u_1),\ldots,y(u_a))|.
\]
Compactness gives a maximum, and spanning gives $D_0>0$. Enumerate rational
tuples $(q_1,\ldots,q_a)$ in $\Omega^a$. Membership and their determinant
$d$ are decided by rational arithmetic. For $d\ne0$, decide the
first-order real formula
\begin{equation}\label{eq:v18-chartformula}
 \forall u_1,\ldots,u_a\in\Omega:\quad
 \det(y(u_1),\ldots,y(u_a))^2<4d^2.
\end{equation}
Quantifier elimination terminates for each formula. Density of rational
points and continuity of determinant imply that some enumerated tuple
has $|d|>D_0/2$. For such a tuple the strict formula is true, so the
search terminates. No computation of an exact algebraic maximizer is
needed.

Invert the rational matrix with columns $y(q_j)$ to obtain $t(q)$.
Cramer's rule and \eqref{eq:v18-chartformula} give $|t_j(q)|<2$ for every
$q$. In the observable norm $\|z\|_* =\max_i|z_i|$, each $v_j$ has norm
at most the diameter of $z(\Omega)$. An upper bound for that diameter can
be produced, for example, by bounding each free coordinate in $[0,1]$
and summing the absolute polynomial coefficients for
$f_i(b(q))-f_i(b(q'))$. This bound is rational and computable. Rational
matrix multiplication and substitution give $x(q)$ and every requested
pullback. They also give its full sparse coefficient encoding. All
searches have now been specified by finite encodings and terminating
operations.
\end{proof}

There is a small but important distinction between the quotient and a
projection onto a smaller set of simulator rows. The former identifies
exactly the directions invisible to the chosen event family; the original
row constraints remain present in every universal positivity test. The
chart need not identify equivalent tasks after a change of marks or of
information signature. Those are changes to the comparison problem.
